\section{Private-pilot estimation and Wald inference}

The efficiency benchmark in \cref{obj:thm:sequential-minimax} is attained by a procedure that uses private pilot releases to select the main-sample channel. The construction combines the pilot treatment-effect estimate with an affine correction computed from the remaining releases. Under interior assignment and positive privacy, a measurable strong-saddle selector exists. The procedure fixes any such selector satisfying the interface in \cref{obj:def:pilot-estimator}: at every interior parameter, it returns feasible staircase weights, a minimizing profiling coordinate, and the optimal information value, with the selected weights also maximizing the objective at that coordinate.

Two restrictions govern the pilot allocation:
\[
m_n\longrightarrow\infty,
\qquad
\frac{m_n}{n}\longrightarrow0,
\]
as specified in \cref{obj:ass:pilot-diverges,obj:ass:pilot-sublinear}. The first supplies an increasing number of private observations for learning the arm means; the second allocates an asymptotically full fraction of subjects to the selected main-sample channel. The attainment result below applies to an arbitrary schedule satisfying these restrictions. The schedule \(m_n=\lfloor\sqrt n\rfloor\) exhibited for the minimax equality in \cref{obj:thm:sequential-minimax} is one such choice.

To describe the estimator, we first specify the selected affine correction \leanref{sym:\widetilde\phi(S)}{\ensuremath{\widetilde\phi(S)}}. It combines the selected profiling direction with the pattern-specific score and normalizes by the selected information value.

\begin{definitionv}{synth_9}[Pilot-selected affine correction]
Let \(p\in(0,1)\), \(\varepsilon>0\), and let \(\widetilde\theta\in(0,1)^2\) be the private pilot estimate. Apply the measurable strong-saddle selector of \cref{obj:def:pilot-estimator} to obtain
\[
(\widetilde\alpha,\widetilde t,\widetilde J)
=(\alpha_{\widetilde\theta},t_{\widetilde\theta},J_{\widetilde\theta}),
\qquad
\widetilde J
=F_{\widetilde\theta}(\widetilde\alpha,\widetilde t)
=J^*(\widetilde\theta,p,\varepsilon)>0.
\]
For each \(S\in\mathcal S\), define the selected affine correction by
\[
\widetilde\phi(S)
=\frac{g_S^{\top}v(\widetilde t)}
       {\widetilde J\,h_S(\widetilde\theta)}.
\]
Here \(v(t)=(t,t+1)^{\top}\), \(h_S(\vartheta)=\pi_\vartheta^{\top}b_S\), and \(g_S=\nabla_\vartheta h_S(\vartheta)\). The selected weights determine the main-sample channel
\[
Q_{\widetilde\alpha}(S\mid x_j)=\widetilde\alpha_S b_S(j)
\]
and the information value
\[
\widetilde J
=\sum_{T\in\mathcal S}
\widetilde\alpha_T
\frac{(g_T^{\top}v(\widetilde t))^2}
     {h_T(\widetilde\theta)}.
\]
For a transcript entry that is a pilot release, set its correction to zero; for a main-sample pattern release \(S\), use \(\widetilde\phi(S)\).
\end{definitionv}

The pilot determines both the release probabilities and the numerical correction attached to each main-sample pattern. Conditional on the private pilot history \(H_{m_n}=(R_1,\ldots,R_{m_n})\), these choices are fixed. The remaining subjects then supply releases under the selected channel. The following boxed algorithm collects the implementation in \cref{obj:def:pilot-estimator}, displaying the steps for sample sizes with \(1\le m_n<n\), as eventually implied by the two pilot-size restrictions.

\begin{center}
\fbox{%
\begin{minipage}{0.94\linewidth}
\textbf{Private-pilot estimation algorithm}

\medskip
\textbf{Inputs.}
The known assignment probability \(p\in(0,1)\), fixed privacy level \(\varepsilon>0\), a pilot schedule satisfying
\[
m_n\to\infty,\qquad m_n/n\to0,
\]
and a fixed measurable strong-saddle selector
\[
\vartheta\longmapsto
(\alpha_\vartheta,t_\vartheta,J_\vartheta)
\]
defined on all real parameter pairs and satisfying all the interior-parameter selection conditions in \cref{obj:def:pilot-estimator}, including maximization over feasible weights at the selected profiling coordinate. The attainment results assume a fixed measurable strong-saddle selector satisfying these conditions. Numerical evaluation of this input requires a separately specified solver and tie-breaking rule.

\begin{enumerate}
\item \textbf{Release and estimate the pilot.}
For the first \(m_n\) subjects, use the four-category \(\varepsilon\)-randomized-response rule in \cref{obj:synth_10}. From the category frequencies, compute
\[
\widehat u_k=\frac{1}{m_n}\sum_{i=1}^{m_n}\mathbf1\{R_i=k\},
\qquad
\widehat\pi_k
=\frac{(e^\varepsilon+3)\widehat u_k-1}{e^\varepsilon-1},
\qquad k\in\{0,1,2,3\}.
\]
With \(\delta_n=(m_n+2)^{-1}\), set
\[
\begin{aligned}
\widetilde\theta_0
&=\max\left\{\delta_n,
\min\left\{1-\delta_n,\frac{\widehat\pi_1}{1-p}\right\}\right\},\\
\widetilde\theta_1
&=\max\left\{\delta_n,
\min\left\{1-\delta_n,\frac{\widehat\pi_3}{p}\right\}\right\}.
\end{aligned}
\]
The zero-based categories \(1\) and \(3\) represent control success and treatment success, respectively.

\item \textbf{Select the main-sample channel.}
Evaluate the input selector at
\(\widetilde\theta=(\widetilde\theta_0,\widetilde\theta_1)^\top\)
to obtain \((\widetilde\alpha,\widetilde t,\widetilde J)\).

\item \textbf{Release the remaining records.}
For \(m_n<i\le n\), generate the pattern release
\[
S_i\sim Q_{\widetilde\alpha}(\,\cdot\mid X_i)
\]
using the selected channel in \cref{obj:synth_9}.

\item \textbf{Estimate the treatment effect.}
Report the pilot contrast plus the average selected correction:
\[
\widehat\tau^*
=\widetilde\theta_1-\widetilde\theta_0
+\frac{1}{n-m_n}\sum_{i=m_n+1}^{n}\widetilde\phi(S_i).
\]
\end{enumerate}
The complete procedure, including its conventions for every sample size, is defined in \cref{obj:def:pilot-estimator}.
\end{minipage}%
}
\end{center}

Clipping places the pilot estimate in the interior parameter space where the selector conditions apply. The main-sample correction then targets the discrepancy between the pilot contrast and the true treatment effect. The conditional unbiasedness clause in \cref{obj:thm:pilot-attainment} gives this centering an exact interpretation: the estimation error has mean zero conditional on the private pilot history. Consequently, the theorem accommodates every pilot schedule satisfying the two stated limits.

For studentization, use the empirical variance estimator \leanref{sym:\widehat V^*}{\ensuremath{\widehat V^*}}. It centers the main-sample corrections at their empirical mean, so the pilot contrast, which is fixed conditional on the pilot history, contributes no term to this empirical variance.

\begin{definitionv}{synth_1}[Empirical variance of selected corrections]
The variance estimator \(\widehat V^*\) is defined for any \(p,\varepsilon\in\mathbb R\), pilot-size schedule \(m:\mathbb N\to\mathbb N\), sample size \(n\in\mathbb N\), transcript \(z\) of \(n\) pilot-or-pattern releases, and selector mapping each real parameter pair to a staircase-weight vector and two real numbers. Write \(m_n=m(n)\), and let \(\widetilde\theta\) be the pilot parameter estimate computed from \(p,\varepsilon,m,z\). For each transcript entry, \(\widetilde\phi(z_i)\) denotes zero when the entry is a pilot release and the selected affine correction \(\widetilde\phi(S_i)\), evaluated at \(\widetilde\theta,p,\varepsilon\) using the supplied selector, when the entry is a pattern release. Then
\[
\widehat V^*
=
\frac{1}{\max\{n-m_n,0\}}
\sum_{\substack{0\le i<n\\m_n\le i}}
\left(
\widetilde\phi(z_i)
-
\frac{1}{\max\{n-m_n,0\}}
\sum_{\substack{0\le j<n\\m_n\le j}}
\widetilde\phi(z_j)
\right)^2.
\]
Here reciprocals are interpreted as zero at zero, so the definition also gives \(\widehat V^*=0\) when \(m_n\ge n\).
\end{definitionv}

For a nonempty main sample, this is the empirical second central moment with divisor \(n-m_n\). Its zero-based transcript indexing corresponds to the one-based main-sample indices in the algorithm. Under the attainment conditions, it consistently estimates the asymptotic variance constant \(V^*(\theta,p,\varepsilon)\) defined in \cref{obj:def:staircase-saddle}. The resulting standard error is \(\sqrt{\widehat V^*/n}\), reflecting the theorem's normalization by the total sample size.

Fix the Wald error probability \leanref{sym:\gamma}{\ensuremath{\gamma}} in \((0,1)\), and let \leanref{sym:\Phi}{\ensuremath{\Phi}} denote the standard normal distribution function. The normal critical value \leanref{sym:z_{1-\gamma/2}}{\ensuremath{z_{1-\gamma/2}}} uses the generalized-inverse convention
\[
z_{1-\gamma/2}
=\inf\{z\in\mathbb R:1-\gamma/2\le\Phi(z)\}.
\]
Together, the selected estimator and empirical variance yield the interval
\[
\left[
\widehat\tau^*-z_{1-\gamma/2}\sqrt{\widehat V^*/n},
\quad
\widehat\tau^*+z_{1-\gamma/2}\sqrt{\widehat V^*/n}
\right].
\]
The next result establishes privacy, conditional centering, regular attainment, variance consistency, and pointwise coverage for this construction.

\begin{theoremv}{thm:pilot-attainment}[Private pilot attainment]
Let \(\theta=(\mu_0,\mu_1)^{\top}\) be the arm-mean parameter, \(p\) the treatment-assignment probability, and \(\varepsilon\) the fixed privacy level. Suppose:
\begin{itemize}
\item \textbf{(Assignment.)} \(p\) satisfies \cref{obj:ass:interior-assignment}.
\item \textbf{(Means.)} \(\theta\) satisfies \cref{obj:ass:interior-means}.
\item \textbf{(Privacy.)} The privacy sequence and \(\varepsilon\) satisfy \cref{obj:ass:fixed-privacy}.
\item \textbf{(Selection.)} The selector in \cref{obj:def:pilot-estimator} satisfies the saddle-selection conditions there and, at every interior parameter \(\vartheta\), its selected weights maximize the objective at its selected profiling coordinate:
\[
F_{\vartheta}(\alpha,t_{\vartheta})
\le F_{\vartheta}(\alpha_{\vartheta},t_{\vartheta})
\qquad
\text{for every }\alpha\in\mathcal A_{\varepsilon}.
\]
\item \textbf{(Pilot schedule.)} The pilot sizes \(m_n\) satisfy \cref{obj:ass:pilot-diverges,obj:ass:pilot-sublinear}.
\item \textbf{(Coverage level.)} The Wald error probability \(\gamma\) satisfies \(0<\gamma<1\).
\end{itemize}

Consider the complete procedure sequence \(\mathcal P\) constructed in \cref{obj:def:pilot-estimator}, with estimator \(\widehat\tau^*\), and write \(\tau=\mu_1-\mu_0\) for the treatment-effect contrast. For every sample size, its channels satisfy sequential \(\varepsilon\)-local differential privacy as specified in \cref{obj:ass:sequential-ldp}.

For \(n\ge2\), let \(H_{m_n}=(R_1,\ldots,R_{m_n})\) be the private pilot history. For every bounded measurable function \(f\) of the transcript whose value depends only on \(H_{m_n}\), the product \((\widehat\tau^*-\tau)f\) is integrable and
\[
\mathbb E_{\theta}\!\left[(\widehat\tau^*-\tau)f\right]=0.
\]
Under every fixed local perturbation \(h\), the scaled local error \(U_{n,h}^{\mathcal P}\), centered at the contrast of the local parameter \(\theta_{n,h}\), satisfies
\[
U_{n,h}^{\mathcal P}
=
\sqrt n\left(
\widehat\tau^*
-
\bigl((\theta_{n,h})_1-(\theta_{n,h})_0\bigr)
\right)
\rightsquigarrow
N\!\left(0,V^*(\theta,p,\varepsilon)\right)
\quad\text{under the transcript law at }\theta_{n,h}.
\]
Here \(V^*(\theta,p,\varepsilon)\) is the optimized contrast variance in \cref{obj:def:staircase-saddle}. In particular,
\[
\sqrt n(\widehat\tau^*-\tau)
\rightsquigarrow N\!\left(0,V^*(\theta,p,\varepsilon)\right).
\]
The variance estimator \(\widehat V^*\) satisfies
\[
\Pr_{\theta}\!\left(
\left|\widehat V^*-V^*(\theta,p,\varepsilon)\right|>\delta
\right)\longrightarrow0
\qquad\text{for every }\delta>0.
\]

The procedure is regular at \(\theta\), with common local limit law
\[
L_{\mathcal P}=N\!\left(0,V^*(\theta,p,\varepsilon)\right).
\]
In addition to the local convergence above, its squared scaled errors satisfy asymptotic uniform integrability over every bounded admissible local-perturbation set: writing \(\mathcal H_{n,H}(\theta)\) for that set at radius \(H\), for every \(H>0\),
\[
\lim_{\substack{M\to\infty\\ M\in\mathbb N}}
\limsup_{n\to\infty}
\sup_{h\in\mathcal H_{n,H}(\theta)}
\mathbb E_{\theta_{n,h}}\!\left[
\bigl(U_{n,h}^{\mathcal P}\bigr)^2
\mathbf 1\!\left\{\left|U_{n,h}^{\mathcal P}\right|>M\right\}
\right]
=0.
\]

Finally, let \(\Phi\) be the standard normal distribution function and define its critical value by
\[
z_{1-\gamma/2}
=
\inf\left\{z\in\mathbb R:1-\gamma/2\le\Phi(z)\right\}.
\]
Then
\[
\Pr_{\theta}\!\left(
\left|\widehat\tau^*-\tau\right|
\le z_{1-\gamma/2}\sqrt{\widehat V^*/n}
\right)\longrightarrow1-\gamma.
\]
\end{theoremv}

The construction separates learning the release rule from centering the treatment-effect estimate. The increasing pilot informs channel selection, while the affine correction exactly centers the estimator conditional on that selection. A vanishing pilot fraction then leaves the main sample to determine the first-order variance. Under the measurable strong-saddle conditions, these features deliver the optimized variance for every admissible pilot schedule, including the exhibited schedule used in \cref{obj:thm:sequential-minimax}.

Regularity strengthens the fixed-parameter normal approximation by giving the same limiting error law under every fixed local perturbation, with centering at the corresponding treatment effect. The additional uniform-integrability clause controls squared-error tails over bounded admissible perturbation sets. Together, these conclusions connect the estimator's distributional behavior to the regular squared-error benchmark in \cref{obj:thm:sequential-minimax}.

The variance-consistency and coverage conclusions apply pointwise at every interior arm-mean parameter satisfying the stated conditions, including parameter points at which the optimal channel support changes. At such points, the measurable strong-saddle selector remains the governing selection rule, and the empirical variance estimates the same optimized variance constant appearing in the normal limit. Thus the procedure supplies a selected private release rule, a treatment-effect estimate, and an asymptotically valid Wald interval throughout this interior domain.
